\section*{第 \ref{ch:b1:electrophilic-additions} 章——烯烃与炔烃：亲电加成}

\begin{solution}{exo:b1:electrophilic-additions:1}
2-氯丙烷；2-溴-2-甲基丙烷；1-溴-1-甲基环己烷；
2-氯丁-2-烯。
\end{solution}

\begin{solution}{exo:b1:electrophilic-additions:2}
1,2-二溴乙烷；\textit{反式}-1,2-二溴环己烷；
1,2-二溴丙烷。
\end{solution}

\begin{solution}{exo:b1:electrophilic-additions:3}
丙-2-醇；2-甲基丙-2-醇；1-甲基环己-1-醇（
每例均符合马氏规则）。
\end{solution}

\begin{solution}{exo:b1:electrophilic-additions:4}
丙炔（\omterm{def:b1:electrophilic-additions:acetylide}{端炔}）：\ce{CH3C#CH + NH2- -> CH3C#C- + NH3}。乙醇：
\ce{CH3CH2OH + NH2- -> CH3CH2O- + NH3}。丁-2-炔和丙烯都没有
酸性足够强的氢。
\end{solution}

\begin{solution}{exo:b1:electrophilic-additions:5}
$\pi$ 电子对夺取 \ce{H+}，加在 \ce{CH2} 碳上，生成叔\omterm{def:b1:electronic-effects:intermediates}{碳正离子}
\ce{(CH3)3C+}；水的\omterm{def:b1:lewis-resonance-vsepr:lewis}{孤电子对}与之成键：\ce{(CH3)3COH2+}；一个水
分子再夺取质子：\ce{(CH3)3COH}。热的浓酸中水很
少，易挥发的烯烃逸出：逆反应，即\omterm{def:b1:alcohol-activation:dehydration}{脱水}，
得以进行。
\end{solution}

\begin{solution}{exo:b1:electrophilic-additions:6}
2-甲基丙烯的质子化生成叔\omterm{def:b1:electronic-effects:intermediates}{碳正离子}，丙烯则生成
仲\omterm{def:b1:electronic-effects:intermediates}{碳正离子}。决速的质子化步骤具有与\omterm{def:b1:electronic-effects:intermediates}{碳正离子}相似的\omterm{def:b1:elementary-steps:energy-profile}{过渡态}
（哈蒙德假设）：较稳定的叔\omterm{def:b1:electronic-effects:intermediates}{碳正离子}
经由较低的能垒生成。
\end{solution}

\begin{solution}{exo:b1:electrophilic-additions:7}
溴正离子桥接环的一面；溴负离子从
另一面开环：两个溴原子最终呈\textit{反式}。\textit{反式}
二溴化物具有\omterm{def:b1:stereochemistry-in-depth:stereocentre}{手性}（无对称面），但溴负离子对两个碳原子的进攻
机会均等：两种\omterm{def:b1:stereochemistry-in-depth:enantiomers}{对映异构体}等量生成，产物为
\omterm{def:b1:stereochemistry-in-depth:enantiomers}{外消旋混合物}，无\omterm{def:b1:stereochemistry-in-depth:optical-activity}{旋光性}。
\end{solution}

\begin{solution}{exo:b1:electrophilic-additions:8}
C1 质子化生成 C2 上的仲\omterm{def:b1:electronic-effects:intermediates}{碳正离子}，与\omterm{def:b1:electronic-effects:carbon-class}{叔碳}
C3（带一个氢）相邻；该氢连同电子对迁至 C2（
1,2-氢负离子迁移），使 C3 上生成叔\omterm{def:b1:electronic-effects:intermediates}{碳正离子}，被
\ce{Cl-} 捕获：2-氯-2-甲基丁烷。
\end{solution}

\begin{solution}{exo:b1:electrophilic-additions:9}
丙炔与氨基钠：生成丙\omterm{def:b1:electrophilic-additions:acetylide}{炔负离子}；与 1-溴丙烷（SN2）作用，
得 \ce{CH3C#C-CH2CH2CH3}，己-2-炔。丙\omterm{def:b1:electrophilic-additions:acetylide}{炔负离子}与苯甲醛作用，
再经稀酸处理，得 \ce{C6H5CH(OH)C#CCH3}，1-苯基丁-2-炔-1-醇。
\end{solution}

\begin{solution}{exo:b1:electrophilic-additions:10}
\cip{E}-丁-2-烯：溴正离子位于一面，溴负离子从另一面进攻 C2
或 C3；两种进攻得到同一化合物，\omterm{def:b1:stereochemistry-in-depth:isomers}{构型}标记为 \cip{2R,3S}：
即\omterm{def:b1:stereochemistry-in-depth:enantiomers}{内消旋化合物}，为一种\omterm{def:b1:stereochemistry-in-depth:isomers}{立体异构体}。\cip{Z}-丁-2-烯：进攻 C2 得
\cip{2S,3S}（或另一面进攻得其镜像），进攻 C3 得
\cip{2R,3R}；两者等量生成，为两种\omterm{def:b1:stereochemistry-in-depth:isomers}{立体异构体}，即\omterm{def:b1:stereochemistry-in-depth:enantiomers}{外消旋混合物}。
\end{solution}

\begin{solution}{exo:b1:electrophilic-additions:11}
在不对称溴正离子中，取代程度较高的碳带有更多的正电荷
（其 C--Br 键更长、更弱，类似近似的\omterm{def:b1:electronic-effects:carbon-class}{叔碳}或仲\omterm{def:b1:electronic-effects:intermediates}{碳正离子}）。水大量过量时，
从与桥相对的一面进攻该碳：
1-溴丙-2-醇。
\end{solution}

\begin{solution}{exo:b1:electrophilic-additions:12}
$D = 1$：烯烃。2-甲基丁-1-烯与 2-甲基丁-2-烯均生成
叔\omterm{def:b1:electronic-effects:intermediates}{碳正离子}，因此分别得 2-溴-2-甲基丁烷与 2-甲基丁-2-醇。二者的
质子核磁共振谱不同：前者有两个乙烯基质子（\ce{=CH2}，单峰），
后者有一个（\ce{=CH-}，四重峰）。
\end{solution}

\begin{solution}{pb:b1:electrophilic-additions:1}
\textbf{1.} 每个碳上 \ce{CH3} 的优先级高于 H：\cip{Z} 中两个甲基
位于同侧，\cip{E} 中位于异侧。
\textbf{2.} 不互为镜像的\omterm{def:b1:stereochemistry-in-depth:isomers}{立体异构体}：\omterm{def:b1:stereochemistry-in-depth:enantiomers}{非对映异构体}。
\textbf{3.} 绕 C=C 键旋转会破坏 $\pi$ 键，
其所需能量远超室温下碰撞所能提供的能量。
\textbf{4.} \cip{E}：其甲基彼此不拥挤。
\textbf{5.} \ce{C4H8 + Br2 -> C4H8Br2}。
\textbf{6.} 两个等价\omterm{def:b1:stereochemistry-in-depth:stereocentre}{手性}碳原子：至多四种，实际上为三种
（\cip{2R,3R}、\cip{2S,3S} 及内消旋的 \cip{2R,3S}）。
\textbf{7.} $\pi$ 电子对进攻 \ce{Br2} 的一个 Br 原子，Br--Br 电子对
以 \ce{Br-} 形式离去，被进攻的溴原子的一个\omterm{def:b1:lewis-resonance-vsepr:lewis}{孤电子对}与
第二个碳成键：生成三元环溴正离子。
\textbf{8.} 桥式离子中每个原子都满足八电子结构；若为开链仲
\omterm{def:b1:electronic-effects:intermediates}{碳正离子}，则该碳只有六个电子。
\textbf{9.} \ce{Br-} 的一个\omterm{def:b1:lewis-resonance-vsepr:lewis}{孤电子对}与一个环碳成键；桥上的 C--Br 键
断裂，其电子对归于桥接的溴原子。
\textbf{10.} 桥占据一个面；与 SN2 反应类似，\omterm{def:b1:electronic-effects:nucleophile}{亲核试剂}从
与将要断裂的键相对的一面进入。
\textbf{11.} \omterm{def:b1:electrophilic-additions:halonium}{反式加成}。
\textbf{12.} 水会与溴负离子竞争溴正离子，
生成溴醇。
\textbf{13.} \cip{2R,3S}（或 \cip{2S,3R}，即同一化合物）。
\textbf{14.} 同一化合物。
\textbf{15.} 在两个溴原子处于\omterm{def:b1:stereochemistry-in-depth:conformations}{重叠式构象}时，一个对称面
通过 C2 与 C3 之间；两者的\omterm{def:b1:stereochemistry-in-depth:isomers}{构型}标记相反。
\textbf{16.} \cip{2R,3R} 与 \cip{2S,3S}。
\textbf{17.} 等量生成：\omterm{def:b1:stereochemistry-in-depth:enantiomers}{外消旋混合物}。
\textbf{18.} 是的：两种非对映的烯烃给出不同产物
（内消旋与外消旋）。
\textbf{19.} 无：两种\omterm{def:b1:stereochemistry-in-depth:enantiomers}{对映异构体}的旋光相互抵消。
\textbf{20.} 各为 $5.60/56.0 = \qty{0.100}{mol}$。
\textbf{21.} 各为 $0.100 \times 0.85 \times 215.8 = \qty{18.3}{g}$。
\textbf{22.} \textbf{$[\alpha] = 0^\circ$（\omterm{def:b1:stereochemistry-in-depth:enantiomers}{内消旋化合物}），\qty{18.3}{g}}。
\end{solution}
